% ---------------------------------------------------------------
\section{Introduction}
\label{sec:introduction}
% ---------------------------------------------------------------

The research system relies on journals and institutions. Journals filter
works through editorial and peer review, archive the literature, andzz
create metadata for retrieval and exploration. Institutions fund individual researchers,
implement training, oversee compliance, establish research organisational units
and maintain physical and information infrastructure. Although the status hierarchy that links journals and institutions is a pervasive aspect of the research system \citep{mertonSociologyScienceTheoretical1973, mcnameeStratificationScienceComparison1994, martinTiedKnowledgePower1998}, bibliometric studies that explore the joint hierarchy of journals and institutions are rare. This paper proposes a joint ordinal spectral ranking of journals and institutions derived from bibliometric data. The method helps to reveal how status is formed and maintained. We also intend to apply the approach to study
the relations between the status of journals and institutions, and researcher esteem.

The introduction of the science citation index \citep{garfieldScienceCitationIndex1964} opened a new way to investigate connections between journals and institutions. Citation-based journal ranking in particular became a useful means to manage collections and to select outlets for publications. On the other hand, the misapplication of the journal impact factor (JIF) as proxy for the importance of research works published in that journal fostered a still-unresolved tension in bibliometric evaluation \citep{garfieldCitationAnalysisLegitimate1979, macrobertsMismeasureScienceCitation2018}. Perhaps due to the complexity of disciplines and missions, an institution impact factor (IIF) has not emerged. Rather, institutional rankings derived from bibliometric indicators of works, authors or journals are often composed with other indicators to form an institutional score. This is the approach in global rankings such as QS, Times Higher Education (THE) and Shanghai Ranking (ARWU), the academic Leiden and SCImago Institutions rankings, as well as many national performance-based research funding schemes \citep{zacharewiczPerformancebasedResearchFunding2019}. 

The prospects for using a citation index to explore research importance were considered in an early work by \cite{narinEvaluativeBibliometricsUse1976}, \cite{pinskiCitationInfluenceJournal1976}, and \cite{narinStructureBiomedicalLiterature1976} who proposed a recursive method to construct influence weights of journals within a disciplinary field. To obtain influence weights for institutions, Narin added publications weighted by the respective journal influence score. Over a sample of several dozen universities, the total institution influence scores were well-correlated with survey-based rankings, butz also with the number of publications, suggesting that survey respondents combined influence and size. Narin's work built on earlier studies of social networks and developed into the eigenvector approach now widely used for journal ranking \citep{bergstrom, vignaSpectralRanking2016, franceschetTenGoodReasons2010}. Eigenvector methods have also been applied to author-level influence weightings in the field of social science \citep{westAuthorlevelEigenfactorMetrics2013}, and extended by Narin's approach to ranking institutions and countries. An alternative approach iterates between the article-article network linked by the reference lists (without projection) and the corresponding article-author network \citep{ujumBestBothWorlds2015, palCITEXNewCitation2015}. The formulation by \cite{caoRankingAcademicInstitutions2023} adopts a similar network, with reference-based article-article links, and the related article-institution links, ranked by a novel and effective iteration over random walks. This is one of very few spectral rankings of institutions. The authors note that a limiting case of this approach corresponds to the eigenvector (pageRank) method. As noted earlier by \cite{gellerCitationInfluenceMethodology1978}, random walks over Markov chains are closely related to the spectral ranking problems discussed in this paper. 

It is useful to foreground analytical devices adopted in this paper. First, we say that a \textit{reference} is an entry $w_2$ in the reference list of a work $w_1$: $w_1$ references $w_2$ and $w_2$ is referenced by $w_1$. Bibliometric indexers establish the connection between $w_2$ and the entry for $w_2$ in an index of the \textit{citations} to $w_2$  \citep{woutersCitationCulture1999}\: we say $w_1$ cites $w_2$ and $w_2$ is cited by $w_1$. A work can be fractionated among its authors and institutions, signifiying the fraction of $w_1$'s \textit{attention} given (paid) to the authors and institutions of $w_2$. Spectral ranking accounts for the weakening and strengthening of $w_1$'s attention by the attention paid to it by others, then distributed by fractionation. We can also say that a citation signifies the \textit{influence} of $w_2$'s authors and institutions on $w_1$; like attention, influence can be weakened or strengthened by other citations to $w_2$. Second, in this paper we view references and citations through Social Systems Citation Theory \citep[][SSCT]{tahamtanSocialSystemsCitation2022}. We concentrate on the structural features of the network of references and citations, setting aside authorial intent. This approach contrasts with actor-network theory \citep[][ANT]{latourReassemblingSocialIntroduction2007}, which foregrounds the intentful roles of 
heterogeneous agents in assembling citation networks. ANT will be important in our follow-up studies that relate this paper to the referencing activities of individual researchers.

This paper proceeds as follows. In the next section, we construct a framework for computing the transfer of attention or influence among journal and institution units coupled by the reference lists of the works they publish. We then construct a corpus of works in the fields of Economics and Business, forming subset of the corpus to illustrate the interconnections of units, fields, and time. We present spectral rankings over these subsets, alongside sensitivity and bootstrap studies of some key parameters. We then discuss the results, including a critique of our methods, and close by mentioning some directions for future work.

